\documentclass[a4paper]{report}
\usepackage[utf8]{inputenc}
\usepackage[portuguese]{babel}
\usepackage{hyperref}
\usepackage{a4wide}
\hypersetup{pdftitle={Sistema de Gestão de Stocks},
pdfauthor={José Ferreira, Emanuel Rodrigues, João Teixeira},
colorlinks=true,
urlcolor=blue,
linkcolor=black}
\usepackage{subcaption}
\usepackage[cache=false]{minted}
\usepackage{listings}
\usepackage{booktabs}
\usepackage{multirow}
\usepackage{appendix}
\usepackage{tikz}
\usepackage{authblk}
\usetikzlibrary{positioning,automata,decorations.markings}

\begin{document}

\title{Gestão de Stocks\\ 
\large Grupo Nº 2}
\author{José Ferreira (A83683) \and Emanuel Rodrigues (A84776) \and João Teixeira (A85504)}
\date{\today}

\begin{center}
    \begin{minipage}{0.75\linewidth}
        \centering
        \includegraphics[width=0.4\textwidth]{eng.jpeg}\par\vspace{1cm}
        \vspace{1.5cm}
        \href{https://www.uminho.pt/PT}
        {\color{black}{\scshape\LARGE Universidade do Minho}} \par
        \vspace{1cm}
        \href{https://www.di.uminho.pt/}
        {\color{black}{\scshape\Large Departamento de Informática}} \par
        \vspace{1.5cm}
        \maketitle
    \end{minipage}
\end{center}

\tableofcontents

\pagebreak

\chapter{Introdução}

O objetivo deste projeto é construir um sistema de gestão de stocks,
de forma a ser capaz de armazenar informação de artigos, stocks e vendas,
aplicando uma arquitetura \textit{Cliente, Servidor} com recurso a \textit{System Calls}
e vários processos, aplicando assim os conhecimentos sobre os Sistema Operativos 
baseados em Unix adquiridos nas aulas desta UC\@.
Como objetivo, é necessário implementar um sistema que permita a utilização concorrente de 
um servidor por vários clientes, assegurando a conservação de toda a informação
introduzida no sistema, bem como a consistência desta entre os vários ficheiros
e cache.\\
Ao longo deste relatório vamos descrever as nossas abordagens a estes problemas.

\chapter{Módulos}\label{chap:api}

\section{Manutenção de Artigos}

O modulo da Manutenção de artigos tem como função ser um `preçário'
tendo informação dos diferentes produtos existentes, bem como o seu
preço. Aqui é possível a criação de novos artigos, bem como a sua
modificação. Sendo um modulo referente a manutenção, também é possível
correr o agregador a pedido.

\section{Servidor de Vendas}

O Servidor de vendas trata todos os pedidos dos Clientes de vendas,
quer sejam vendas ou entradas de stock, ou informações sobre o preço
e stocks existentes no momento.

\section{Cliente de Vendas}

O Cliente de Vendas é a interface que o utilizador final usa para
realizar todas as vendas e pedidos de informação, comunicando com
o servidor de vendas.

\section{Agregador}

O Agregador tem como finalidade agregar vendas, ou seja, juntar
todas as vendas referentes ao mesmo artigo.

\chapter{Arquitetura e Solução do Projeto}

\section{Manutenção de Vendas}

A Manutenção de Vendas trata de manter artigos, respetivos preços
e nomes. Para manter essa informação usa dois ficheiros, \textit{strings}
e \textit{artigos}. O utilizador consegue adicionar e modificar artigos
através do stdin. Embora a manutenção de vendas seja independente dos 
restantes módulos, este comunica com o Servidor de Vendas, caso este esteja
a correr, de forma a manter a sincronização  da informação sobre o preço 
e numero de artigos existentes em ambos os lados.\\
Aqui também é possível correr o agregador sobre o ficheiro de vendas.

\section{Servidor de Vendas}

O Servidor de Vendas trata dos pedidos de todos os clientes, e trata de 
manter informação sobre o stock de cada artigo. Para manter essa informação
utiliza o ficheiro \textit{stocks}. O Servidor corre em background, consistindo
em três diferentes processos:
\begin{itemize} 
    \item Um que trata da sincronização com a
        parte da Manutenção de Artigos (preços e criação de novos artigos), com a
        finalidade de manter, tanto o ficheiro de stocks, como a cache de preços
        sempre atualizada.
    \item Um que contém a cache de preços.
    \item Um responsável pela comunicação e tratamento dos pedidos dos clientes.
\end{itemize}
Toda a comunicação com o servidor é feita com recurso a \textit{Named pipes},
existindo um pipe para os clientes efetuarem pedidos e um para o Manutenção
de Artigos enviar as alterações que ocorram da parte dele.

\section{Cliente de Vendas}

O Cliente de Vendas interage com o utilizador a partir do \textit{stdin} e 
\textit{stdout}, e interage com o Servidor de Vendas enviando informação
para o pipe de entrada que o Servidor cria, e recebe a resposta deste
através de um pipe que ele cria, tendo como identificador o seu \textit{PID},
de forma a garantir que as informações não são lidas por outro cliente.
Para isto acontecer, para além do pedido, e escrito no pipe de leitura do
servidor o identificador do cliente, que este esta preparado para conseguir
responder no pipe adequado.

\section{Agregador}

O agregador recebe input do \textit{stdin} linhas com o formato estipulado
para o ficheiro \textit{vendas}, calculando os totais de cada artigo,
devolvendo esses totais no mesmo formato pelo \textit{stdout}.

\chapter{Estruturas de dados}

Como descrito no capítulo\ref{chap:api}, todas as nossas estruturas de dados consistem 
em Hashtables da GLib, escolha feita após múltiplos testes empíricos de abordagens distintas.\\
Numa abordagem inicial, na primeira fase do projeto, utilizamos arrays para 
guardar os dados. Para esta abordagem, observamos tempos de leitura e validação de ficheiros
superiores a 5 minutos.\\
Numa segunda abordagem, no inicio da segunda fase, optamos por utilizar 
Árvores Auto Balanceadas da biblioteca GLib, onde os tempos de leitura e validação 
dos ficheiros baixaram para cerca de 7 segundos, graças à pesquisa $O(\log N)$,
operação mais utilizada na validação das vendas.\\
Como abordagem final, implementamos Hashtables, também da biblioteca GLib.
Nesta, os valores não são armazenados de forma ordenada. No entanto,
tal como concluímos com esta versão, com ferramentas tão otimizadas para a ordenação
como o \textit{qsort()} da linguagem de programação C, a inserção ordenada acaba por ter
um impacto demasiado negativo nos tempos de inserção.\\
Por um lado, graças ao tempo de inserção e pesquisa $O(1)$, conseguimos reduzir
o tempo de carregamento de ficheiros para menos de 4 segundos, como descrito
na tabela, um melhoramento de aproximadamente 40\%.\\
Por outro lado, graças à implementação otimizada das Hashtables na biblioteca em uso,
percorrer todos os elementos desta, de acordo com os nossos testes, é tão rápido como
na implementação anterior de AVLs. Tal resulta num tempo de resposta a queries
inalterado quando comprado com a implementação anterior.\\
Assim, com a utilização de Hashtables conseguimos uma melhoria de
aproximadamente 40\% nos tempos de \textit{load} e manter os mesmos resultados nos
tempos de resposta a queries.

\chapter{Testes e Benchmarks}

Durante a execução do nosso projeto, foram efetuados diversos testes ao código, com
ferramentas como \textit{Valgrind}, para verificar a presença de \textit{memory 
leaks}, \textit{memusage}, para fazer profiling à memoria utilizada, \textit{gprof}
para ver quais funções pesam mais na execução do programa, com vista a otimizar as
partes mais lentas, e para medição dos tempos de execução utilizamos a biblioteca
\textit{time.h}.

\section{Memória}

Utilizando a ferramenta \textit{Valgrind}, conseguimos verificar que o programa
corre sem \textit{memory leaks}, à exceção de algumas causadas pela biblioteca
\textit{glib}.\\

Com a ferramenta \textit{memusage}, obtivemos o gráfico de uso de memória presente
na figura, concluindo que as nossas estruturas preenchidas com
os dados dados, estão a consumir, em pico, por volta dos 360MB.

\chapter{Conclusão}

Para concluir, conseguimos cumprir todos os requisitos propostos, conseguindo implementar
todos os módulos e estruturá-los como pedido, sendo assim capaz de responder a todas as 
queries da forma que nos pareceu mais eficiente.\\
Como trabalho futuro, gostaríamos de conseguir melhorar os tempos de load de ficheiros,
no caso do ficheiro de um milhão de vendas, para tempos inferiores a 2 segundos, e
gostaríamos também de fazer ligeiras mudanças à estruturação das APIs de alguns dos 
módulos, para no futuro, ser mais fácil de responder a novas queries caso fosse 
necessário.

\appendix

\end{document}
